\documentclass[11pt,a4paper]{article}
\usepackage[utf8]{inputenc}
\usepackage{amsmath,amssymb,amsfonts,amsthm}
\usepackage{booktabs}
\usepackage{geometry}
\usepackage{hyperref}
\usepackage{algorithm}
\usepackage{algorithmic}
\usepackage{graphicx}
\usepackage{xcolor}
\usepackage{cite}
\geometry{margin=1in}

\title{\textbf{Cognitively Constrained Heritage Trails \& Route Optimization: \\
Topological Route Planning with Epistemic Verification in Cultural Field Inquiry (LN-06)}}
\author{\textbf{Lưu Niên Sứ Ký Research Consortium} \\
Huế Monuments Conservation Center \& Institute of Information Technology}
\date{November 2026}

\newtheorem{definition}{Definition}
\newtheorem{theorem}{Theorem}
\newtheorem{invariant}{Invariant}
\newtheorem{ruleitem}{Rule}

\begin{document}
\maketitle

\begin{abstract}
Designing field-based learning itineraries across complex cultural heritage landscapes---such as the Imperial Citadel of Huế and the UNESCO Memory of the World Nine Dynastic Urns (\textit{Cửu Đỉnh})---presents a fundamental tension between spatial proximity and cognitive scaffolding. Conventional tourist routing systems employ spatial Traveling Salesperson Problem (TSP) heuristics that minimize geographic transit distance, routinely leading to severe pedagogical incoherence where advanced synthesis tasks precede foundational epistemological encounters. In this paper, we formalize the \textit{Cognitively Constrained Heritage Trail} as a directed acyclic graph (DAG) routing problem governed by five domain invariants: prerequisite dependency satisfaction, source authenticity verification, strict temporal budgeting, outcome coverage, and acyclicity. We implement an exact polynomial-time topological planning engine and evaluate it against an independent 120-case holdout dataset spanning five declared fault archetypes (\textit{nominal}, \textit{unreachable\_checkpoint}, \textit{source\_gap\_in\_monument}, \textit{prerequisite\_cycle}, and \textit{over\_time\_budget}). While standard spatial TSP heuristics ($B_0$) and greedy topological ordering without budget bounding ($B_1$) achieve only $20.00\%$ accuracy with up to $2.80$ prerequisite violations per route ($336$ total), our cognitively constrained policy ($I_1$) achieves $100.00\%$ decision accuracy [$95\%$ CI: $96.90\%$--$100.00\%$], zero prerequisite violations ($p < 10^{-21}$), and rigorous rejection of uncorroborated heritage stops. Ablation experiments ($A_1$--$A_3$) demonstrate that isolating any individual constraint introduces deterministic pedagogical failures, validating the indivisible nature of cognitive, temporal, and epistemic bounds in educational heritage field studies.
\end{abstract}

\section{Introduction}
Spatial orientation in open-air cultural monuments is rarely aligned with instructional progression. In the Imperial City of Huế (Đại Nội), physical monuments are arranged according to classical geomantic and dynastic axes: entering through the South Meridian Gate (\textit{Ngọ Môn}), traversing the imperial court to the Palace of Supreme Harmony (\textit{Điện Thái Hòa}), turning westward to the Ancestral Temple (\textit{Thế Tổ Miếu}), and accessing the courtyard sheltering the Nine Dynastic Urns (\textit{Cửu Đỉnh}). 

When automated mobile tour apps or digital learning management systems plan student field trips, they universally treat itinerary generation as a spatial optimization problem:
\begin{equation}
\min_{\pi} \sum_{i=1}^{n-1} d(s_{\pi(i)}, s_{\pi(i+1)})
\end{equation}
where $d(\cdot, \cdot)$ denotes physical Euclidean or walking network distance. However, cultural monuments encode dense hierarchical meaning. For example, understanding the bas-relief motifs of dynastic sovereignty on the bronze urns (e.g., the Eastern Sea \textit{Biển Đông} on Cao Đỉnh or the Bạch Đằng River on Huyền Đỉnh) requires prior mastery of:
\begin{enumerate}
    \item Physical and spatial orientation within the royal precinct (S1 competency);
    \item Dynastic chronology and ceremonial rank under Emperor Minh Mạng (S2 competency);
    \item Primary source textual citations from the \textit{Đại Nam Thực Lục} (S3 competency);
    \item Multi-source historical argumentation and territorial sovereignty claims (S4 competency).
\end{enumerate}

When unconstrained spatial algorithms schedule checkpoints purely based on proximity, learners are repeatedly assigned high-order analytical tasks before acquiring foundational contextual knowledge. In empirical heritage education, this inversion results in severe cognitive overload, fragmented recall, and reliance on unverified internet speculation.

This investigation introduces Study \textbf{LN-06}: a formal framework for \textit{Cognitively Constrained Heritage Trails}. We formulate trail generation as a multi-constrained topological sort over a pedagogical dependency graph $(V, E)$, bounded by real-world educational field constraints: 45-minute or 90-minute secondary school class period allocations, physical checkpoint reachability, and cryptographic primary-source authentication.

\section{Formal Problem Formulation \& Domain Invariants}

\subsection{Mathematical Formulation}
Let a heritage site be represented by a set of candidate stops $S = \{s_1, s_2, \dots, s_m\}$. Each stop $s_i$ is characterized by a tuple:
\begin{equation}
s_i = \langle \text{id}_i, c_i, P_i, R_i, v_i \rangle
\end{equation}
where:
\begin{itemize}
    \item $c_i \in \mathbb{Z}^+$ is the dwell and transit cost in minutes;
    \item $P_i \subseteq \mathcal{O}$ is the finite set of pedagogical learning outcomes provided upon completing checkpoint $s_i$;
    \item $R_i \subseteq \mathcal{O}$ is the set of required prerequisite outcomes that must be acquired prior to visiting $s_i$;
    \item $v_i \in \{0, 1\}$ is an epistemic verification indicator signifying whether the historical assertions linked to $s_i$ are grounded in primary historical records (\textit{Châu bản triều Nguyễn}, \textit{Đại Nam Thực Lục}).
\end{itemize}

A field inquiry objective is defined by $\mathcal{Q} = \langle \mathcal{O}^*, B, A \rangle$, where $\mathcal{O}^* \subseteq \mathcal{O}$ represents target competency outcomes, $B \in \mathbb{Z}^+$ is the total class time budget in minutes, and $A \subseteq S$ denotes the subset of currently accessible physical monuments (accounting for site restorations, weather closures, or ticketing boundaries).

\subsection{Core Domain Invariants}
The trail planning engine enforces five foundational invariants derived from the Vietnam 2018 General Education Program (GDPT 2018) for History and Heritage:

\begin{invariant}[\textbf{INV-LN-06-01: Prerequisite Precedence}]
For any planned route sequence $\pi = (s_{\pi(1)}, s_{\pi(2)}, \dots, s_{\pi(k)})$, no stop may be scheduled unless all its prerequisite competencies have been provided by preceding stops:
\begin{equation}
\forall j \in \{1, \dots, k\}, \quad R_{\pi(j)} \subseteq \bigcup_{i=1}^{j-1} P_{\pi(i)}
\end{equation}
\end{invariant}

\begin{invariant}[\textbf{INV-LN-06-02: Temporal Budget Feasibility}]
The total temporal cost of the itinerary cannot exceed the scheduled pedagogical class budget $B$:
\begin{equation}
\text{Cost}(\pi) = \sum_{j=1}^{k} c_{\pi(j)} \le B
\end{equation}
\end{invariant}

\begin{invariant}[\textbf{INV-LN-06-03: Reachability Completeness}]
A trail is feasible if and only if all stops required to achieve target outcomes $\mathcal{O}^*$ belong to the set of physically accessible stops $A$:
\begin{equation}
S_{\text{needed}}(\mathcal{O}^*) \subseteq A
\end{equation}
\end{invariant}

\begin{invariant}[\textbf{INV-LN-06-04: Epistemic Source Gating}]
No learning itinerary may include stops whose historical claims lack archival corroboration:
\begin{equation}
\forall s_i \in \pi, \quad v_i = 1
\end{equation}
If any candidate checkpoint exhibits $v_i = 0$, the system must emit \texttt{reject\_checkpoint\_source}.
\end{invariant}

\begin{invariant}[\textbf{INV-LN-06-05: Acyclicity of Cognitive Dependencies}]
The dependency graph $G = (S, E)$ where $(u, v) \in E \iff P_u \cap R_v \ne \emptyset$ must be a Directed Acyclic Graph (DAG). If cycles exist ($G$ contains a directed cycle), the engine must fail-closed with \texttt{reject\_cycle}.
\end{invariant}

\section{Algorithmic Implementation}
The intervention policy $I_1$ executes the five-stage evaluation detailed in Algorithm~\ref{alg:cognitive_planner}.

\begin{algorithm}[ht]
\caption{Cognitively Constrained Heritage Trail Planner ($I_1$)}
\label{alg:cognitive_planner}
\begin{algorithmic}[1]
\REQUIRE Stops $S$, Available $A$, Target Outcomes $\mathcal{O}^*$, Budget $B$, Verified Flag $V$
\ENSURE RoutePlan $\langle \text{decision}, \pi, \text{cost}, \text{coverage}, \text{violations} \rangle$
\IF{$\neg V$}
    \RETURN $\langle \texttt{"reject\_checkpoint\_source"}, [], 0, 0.0, 0 \rangle$
\ENDIF
\STATE Construct dependency graph $G = (S, E)$ from $P_u$ and $R_v$
\IF{ContainsCycle($G$)}
    \RETURN $\langle \texttt{"reject\_cycle"}, [], 0, 0.0, 0 \rangle$
\ENDIF
\STATE $S_{\text{needed}} \leftarrow \text{TransitiveClosure}(S, \mathcal{O}^*)$
\IF{$S_{\text{needed}} \not\subseteq A$ \textbf{or} $\mathcal{O}^* \not\subseteq \bigcup_{s \in A} P_s$}
    \STATE $r_{\text{unreach}} \leftarrow |S_{\text{needed}} \setminus A| / |S_{\text{needed}}|$
    \RETURN $\langle \texttt{"no\_feasible\_route"}, [], 0, \text{cov}, 0, r_{\text{unreach}} \rangle$
\ENDIF
\STATE $\text{TotalCost} \leftarrow \sum_{s \in S_{\text{needed}}} c_s$
\IF{$\text{TotalCost} > B$}
    \RETURN $\langle \texttt{"cost\_exceeded"}, [], \text{TotalCost}, 1.0, 0 \rangle$
\ENDIF
\STATE $\pi \leftarrow \text{KahnTopologicalSort}(S_{\text{needed}}, E)$
\STATE $\text{violations} \leftarrow \text{ComputePrerequisiteViolations}(\pi)$
\RETURN $\langle \texttt{"find\_valid\_route"}, \pi, \text{TotalCost}, 1.0, \text{violations} \rangle$
\end{algorithmic}
\end{algorithm}

\section{Experimental Methodology \& Independent Holdout}

\subsection{Holdout Dataset Design}
To evaluate model generalization beyond training fixtures, we constructed an independent 120-case holdout dataset:
\begin{equation}
\mathcal{D}_{\text{holdout}} = \{ (x_i, y_i^*, \text{expert}_i) \}_{i=1}^{120}
\end{equation}
stored at \texttt{research/fixtures/holdout/LN-06\_holdout.jsonl}. The holdout comprises exactly five balanced fault archetypes ($n = 24$ cases each):
\begin{enumerate}
    \item \textbf{Nominal ($n=24$):} Valid DAG configurations with sufficient budgets ($B = 75$ min), full physical accessibility, and verified primary sources.
    \item \textbf{Unreachable Checkpoint ($n=24$):} Intermediate scaffolding stops removed from $A$ ($|S_{\text{needed}} \setminus A| \ge 1$), testing whether planners naively skip steps or accurately detect pedagogical breakages.
    \item \textbf{Source Gap in Monument ($n=24$):} Checkpoints linked to unreviewed secondary assertions ($v_i = 0$), testing epistemic integrity gates.
    \item \textbf{Prerequisite Cycle ($n=24$):} Circular cognitive dependencies introduced (e.g., $s_1$ requires outcome of $s_4$ which requires $s_1$), testing cycle detection.
    \item \textbf{Over Time Budget ($n=24$):} Tight classroom constraints ($B = 30$ min, route cost $\approx 70$ min), testing cognitive load and duration enforcement.
\end{enumerate}

\subsection{Comparative Conditions}
We evaluate six experimental conditions:
\begin{itemize}
    \item \textbf{$B_0$ (Naive Spatial TSP):} Visits available stops in geographic traversal order (modeled as outer-to-inner or reverse layout), ignoring prerequisite dependencies, budget caps, and source checks.
    \item \textbf{$B_1$ (Greedy Topological Search):} Attempts topological sort but lacks budget constraints, permits cycles by blindly appending remaining nodes, and lacks source verification.
    \item \textbf{$I_1$ (Cognitively Constrained DAG Planner):} The proposed full intervention enforcing all five domain invariants.
    \item \textbf{$A_1$ (Ablation 1 - Disable Prerequisites):} Deactivates cycle checks and prerequisite topological sequencing.
    \item \textbf{$A_2$ (Ablation 2 - Disable Budget):} Deactivates class duration upper bounds.
    \item \textbf{$A_3$ (Ablation 3 - Disable Sources):} Deactivates source verification gating.
\end{itemize}

\section{Empirical Results}

\subsection{Overall Benchmark Performance}
Table~\ref{tab:holdout_results} summarizes the performance across all 120 independent holdout scenarios.

\begin{table*}[t]
\centering
\small
\caption{Empirical evaluation of trail planning policies on the 120-case independent holdout dataset.}
\label{tab:holdout_results}
\begin{tabular}{llcccccc}
\toprule
\textbf{Cond.} & \textbf{Policy Description} & \textbf{Acc. (\%)} & \textbf{95\% CI (Wilson)} & \textbf{Mean Viol.} & \textbf{Total Viol.} & \textbf{Outcome Cov.} & \textbf{Mean Cost (m)} \\
\midrule
$B_0$ & Naive Spatial TSP & 20.00\% & [13.82\%, 28.04\%] & 2.800 & 336 & 0.950 & 67.50 \\
$B_1$ & Greedy Topological (No Budget) & 20.00\% & [13.82\%, 28.04\%] & 0.400 & 48 & 0.950 & 67.50 \\
\textbf{$I_1$} & \textbf{Cognitive DAG Planner (Ours)} & \textbf{100.00\%} & \textbf{[96.90\%, 100.00\%]} & \textbf{0.000} & \textbf{0} & \textbf{0.550} & \textbf{28.50} \\
$A_1$ & Ablation: Disable Prerequisites & 80.00\% & [71.96\%, 86.18\%] & 0.200 & 24 & 0.750 & 42.75 \\
$A_2$ & Ablation: Disable Budget & 80.00\% & [71.96\%, 86.18\%] & 0.000 & 0 & 0.550 & 28.50 \\
$A_3$ & Ablation: Disable Sources & 80.00\% & [71.96\%, 86.18\%] & 0.000 & 0 & 0.750 & 42.75 \\
\bottomrule
\end{tabular}
\end{table*}

\subsection{Fault Type Diagnostic Breakdown}
Table~\ref{tab:fault_breakdown} delineates the itemized accuracy across the five fault archetypes.

\begin{table}[ht]
\centering
\small
\caption{Accuracy breakdown across declared fault archetypes ($n=24$ cases per fault).}
\label{tab:fault_breakdown}
\begin{tabular}{lcccccc}
\toprule
\textbf{Fault Archetype} & \textbf{$B_0$} & \textbf{$B_1$} & \textbf{$I_1$} & \textbf{$A_1$} & \textbf{$A_2$} & \textbf{$A_3$} \\
\midrule
Nominal ($n=24$) & 100\% & 100\% & 100\% & 100\% & 100\% & 100\% \\
Unreachable Checkpoint ($n=24$) & 0\% & 0\% & 100\% & 100\% & 100\% & 100\% \\
Source Gap in Monument ($n=24$) & 0\% & 0\% & 100\% & 100\% & 100\% & 0\% \\
Prerequisite Cycle ($n=24$) & 0\% & 0\% & 100\% & 0\% & 100\% & 100\% \\
Over Time Budget ($n=24$) & 0\% & 0\% & 100\% & 100\% & 0\% & 100\% \\
\midrule
\textbf{Overall ($N=120$)} & \textbf{20\%} & \textbf{20\%} & \textbf{100\%} & \textbf{80\%} & \textbf{80\%} & \textbf{80\%} \\
\bottomrule
\end{tabular}
\end{table}

\subsection{Statistical Significance \& Hypothesis Testing}
To determine whether the observed performance gains of $I_1$ over baselines and ablations are statistically significant, we conducted:
\begin{enumerate}
    \item McNemar's test with continuity correction on binary decision correctness ($df = 1$);
    \item Paired two-tailed Student's $t$-tests on the per-route prerequisite violation count.
\end{enumerate}

As presented in Table~\ref{tab:significance}, $I_1$ achieves decisive statistical superiority over both baselines ($p < 10^{-21}$ on McNemar; $p < 10^{-18}$ on paired violations). Every single ablation condition ($A_1$--$A_3$) exhibits an exact 24-case failure mode corresponding directly to its removed invariant ($p = 2.67 \times 10^{-6}$).

\begin{table}[ht]
\centering
\small
\caption{Statistical significance comparisons against Intervention $I_1$.}
\label{tab:significance}
\begin{tabular}{lcccc}
\toprule
\textbf{Comparison} & \textbf{McNemar $\chi^2$} & \textbf{McNemar $p$} & \textbf{Paired $t$ (Viol.)} & \textbf{$t$-test $p$} \\
\midrule
$I_1$ vs. $B_0$ & 94.01 & $3.14 \times 10^{-22}$ & 76.36 & $< 10^{-30}$ \\
$I_1$ vs. $B_1$ & 94.01 & $3.14 \times 10^{-22}$ & 8.91 & $5.25 \times 10^{-19}$ \\
$I_1$ vs. $A_1$ & 22.04 & $2.67 \times 10^{-6}$ & 5.45 & $4.92 \times 10^{-8}$ \\
$I_1$ vs. $A_2$ & 22.04 & $2.67 \times 10^{-6}$ & 0.00 & $1.00$ \\
$I_1$ vs. $A_3$ & 22.04 & $2.67 \times 10^{-6}$ & 0.00 & $1.00$ \\
\bottomrule
\end{tabular}
\end{table}

\section{Discussion \& Pedagogical Implications}

\subsection{The Illusion of High Outcome Coverage}
An intriguing phenomenon observed in Table~\ref{tab:holdout_results} is that baselines $B_0$ and $B_1$ report higher apparent \textit{Outcome Coverage} ($0.950$) than the intervention policy $I_1$ ($0.550$). In superficial LMS metrics, $B_0$ might appear superior. However, this is an artifact of fail-open scheduling:
\begin{itemize}
    \item When a route is physically infeasible (missing intermediate stops), contains cycles, or exceeds time budgets, $B_0$ and $B_1$ naively generate an unfeasible route and tally all visited stops.
    \item In contrast, $I_1$ correctly rejects invalid routes, yielding $0.0$ coverage on failed trails while achieving $1.000$ coverage on all nominal routes.
\end{itemize}
This underscores the vital necessity of evaluating policy decisions against formal ground-truth oracles rather than unconstrained reward metrics.

\subsection{Field Study Constraints in Huế Heritage Education}
During secondary school field excursions in Huế, teachers manage cohorts of 30--40 students under tight 45-minute or 90-minute periods. An unconstrained trail that overruns the schedule by 25 minutes ($67.5$ min vs. $30.0$ min budget in $B_0$) disrupts school bus logistics, forces instructors to rush through the Nine Urns without substantive source observation, and degrades historical inquiry into superficial sightseeing. 

Furthermore, by strictly enforcing \textbf{INV-LN-06-04}, our engine prevents the proliferation of historical misconceptions. During the colonial era and early modern tourism, several motifs on Cửu Đỉnh were erroneously romanticized or misattributed in folk guides. By grounding every checkpoint in verified epigraphy from \textit{Đại Nam Thực Lục Chính Biên}, the platform guarantees that student inquiry is anchored in authentic archival records.

\section{Conclusion}
Study LN-06 demonstrates that heritage trail generation cannot be treated as a purely spatial distance minimization problem. By formulating trail construction as a multi-constrained DAG optimization problem over pedagogical dependencies, archival sources, and temporal budgets, the Lưu Niên platform eliminates prerequisite violations, prevents instructional deadlocks, and guarantees epistemic fidelity. Benchmark evaluation on an independent 120-case holdout dataset confirms that the intervention achieves $100\%$ policy correctness, whereas conventional spatial and unconstrained topological heuristics fail in $80\%$ of non-nominal pedagogical scenarios.

\bibliographystyle{IEEEtran}
\begin{thebibliography}{10}

\bibitem{nguyen2026cuudinh}
Huế Monuments Conservation Center, \emph{The Nine Dynastic Urns of Huế: UNESCO Documentary Heritage}, Thuận Hóa Publishing House, Huế, 2024.

\bibitem{moet2018gdpt}
Ministry of Education and Training (MOET) Vietnam, \emph{General Education Program: History Curriculum (Grades 10--12)}, Hanoi, 2018.

\bibitem{kahn1962topological}
A.~B. Kahn, ``Topological sorting of large networks,'' \emph{Communications of the ACM}, vol.~5, no.~11, pp. 558--562, 1962.

\bibitem{zimmerman2000self}
B.~J. Zimmerman, ``Attaining self-regulation: A social cognitive perspective,'' in \emph{Handbook of Self-Regulation}, Academic Press, pp. 13--39, 2000.

\bibitem{vanlehn2006behavior}
K.~VanLehn, ``The behavior of tutoring systems,'' \emph{International Journal of Artificial Intelligence in Education}, vol.~16, no.~3, pp. 227--265, 2006.

\end{thebibliography}

\end{document}
